\chapter{Ajustements de financement, de marge et de capital}\label{ch:rc:funding-margin-and-capital-adjustments}

En 2012, John Hull et Alan White se sont demandé par écrit si le financement constitue un coût pour les
desks de produits dérivés, et ont soutenu qu'il ne devait pas entrer dans le prix d'un dérivé. Une réponse de praticien
parut la même année sous le titre «~Oui, la FVA est un coût pour les desks de produits dérivés~». En janvier
2014, JPMorgan annonça une perte de 1.5~milliard de dollars liée à la mise en œuvre d'un \omterm{def:rc:funding-margin-and-capital-adjustments:fva}{ajustement de valeur pour financement},
décrivant une migration de l'industrie vers la tarification du coût de financement. Les deux camps avaient raison sur un point~: le débat portait sur la
valeur à mesurer, celle de la firme dans son ensemble ou celle de ses actionnaires, et sur ce
qu'une transaction non collatéralisée oblige une banque à emprunter. Ce chapitre construit les ajustements qui
ont suivi la CVA et la DVA~: le coût de financement des positions non collatéralisées, celui de la marge initiale
à verser, et celui du capital détenu au titre du risque de contrepartie et du risque de CVA, puis les empile dans le
prix que la banque affiche.

\section{Financer un dérivé~: le débat}

Une banque qui conclut avec un client un swap non collatéralisé le couvre au moyen d'un swap opposé
conclu sous un CSA avec le marché. Lorsque le swap du client est favorable à la banque, la
couverture joue contre elle~: la banque verse une marge de variation sur la couverture et ne reçoit rien du
client. Elle doit emprunter ces liquidités à son taux de financement propre, au-dessus du taux overnight versé sur le
collatéral. Lorsque le swap du client est défavorable à la banque, celle-ci reçoit la marge sur
la couverture et peut la prêter ou réduire son endettement
(voir la \cref{fig:rc:funding-margin-and-capital-adjustments:flows}).

\begin{omfigure}
\begin{tikzpicture}[font=\small,box/.style={draw,rounded corners,minimum width=2.4cm,minimum height=0.9cm,align=center}]
\node[box] (c) at (0,0) {client\\(sans CSA)};
\node[box,fill=omDef!10] (b) at (5,0) {banque};
\node[box] (h) at (10.4,0) {contrepartie\\de couverture (CSA)};
\node[box,fill=omThm!10] (t) at (5,-2.8) {trésorerie de la banque\\(emprunte à $r+s_f$)};
\draw[<->,thick] (c) -- node[above] {swap} (b);
\draw[<->,thick] (b) -- node[above] {swap opposé} (h);
\draw[->,omThm,thick] (b.south east) to[bend right=25] node[below,align=center] {marge de variation versée\\quand le client doit\\à la banque} (h.south west);
\draw[->,omThm,thick] (t) -- node[left,align=right] {liquidités\\empruntées} (b);
\end{tikzpicture}

\omcaption{Pourquoi une transaction non collatéralisée requiert un financement~: la banque la couvre au moyen d'une
transaction collatéralisée, verse une marge sur la couverture lorsque la transaction du client lui est favorable,
et emprunte cette marge à sa trésorerie au spread de financement qui lui est propre. Schéma indicatif.}\label{fig:rc:funding-margin-and-capital-adjustments:flows}
\end{omfigure}

\begin{definition}[Ajustement de valeur pour financement]\label{def:rc:funding-margin-and-capital-adjustments:fva}
L'\emph{ajustement de valeur pour financement}\index{ajustement de valeur pour financement} (funding valuation adjustment, FVA) est la
valeur du coût de financement du collatéral que requièrent les transactions non collatéralisées et
leurs couvertures, diminuée du bénéfice du collatéral qu'elles libèrent~:
$\mathrm{FVA} = \mathrm{FCA}-\mathrm{FBA}$.
\end{definition}

\begin{definition}[Ajustements pour coût et bénéfice de financement]\label{def:rc:funding-margin-and-capital-adjustments:fca}
L'\emph{ajustement pour coût de financement}\index{ajustement pour coût de financement} (funding cost adjustment) est
$\mathrm{FCA} = \int_0^Ts_f(t)\,\mathrm{DEE}(t)\,Q_C(t)Q_B(t)\,dt$, le spread de financement $s_f$ versé sur
l'\omterm{def:rc:counterparty-exposure:profiles}{exposition positive attendue} actualisée tant que les deux parties survivent~; le \emph{ajustement pour
bénéfice de financement}\index{ajustement pour bénéfice de financement} (FBA) est la même intégrale appliquée à
l'\omterm{def:rc:counterparty-exposure:profiles}{exposition négative attendue} actualisée, au taux qu'économise la banque.
\end{definition}

La FBA recoupe la DVA~: les deux représentent une valeur que la banque capte grâce à son propre spread de crédit
sur ce qu'elle doit. Selon les arguments de valeur de la firme (Hull et White), la FVA ne devrait pas affecter les prix~;
en revanche, les arguments de valeur pour l'actionnaire (Andersen, Duffie et Song) font valoir que les coûts de financement transfèrent
de la valeur aux créanciers, et qu'un desk doit les répercuter. Les conventions diffèrent d'une banque à l'autre~: certaines comptent
la FCA et la FBA de façon symétrique, d'autres seulement le coût, d'autres encore neutralisent le chevauchement avec la DVA.

\section{Ajustements pour coût et bénéfice de financement}

\begin{example}[FVA du swap à dix ans]\label{ex:rc:funding-margin-and-capital-adjustments:fva}
Prenons le swap à dix ans de 100~millions de dollars du chapitre~17, sur lequel la banque reçoit le taux
au pair de 3.80~\% auprès de la contrepartie BBB du chapitre~18, et une banque se finançant à 80 points de base
au-dessus du taux overnight. Les deux parties étant supposées survivre, la FCA s'élève à 127\,400~dollars et la
FBA de 185\,797~dollars~: le swap doit être attendu comme un passif plutôt qu'un actif, si bien que la FVA nette
est un bénéfice de 58\,397~dollars. La FCA seule, une charge que beaucoup de banques comptabilisent, vaut 1.5~points de base par
an sur le swap. Sous un CSA avec marge quotidienne, l'exposition de financement presque disparaît~: la
FCA tombe à 3\,102~dollars.
\end{example}

\omcode{code/firm/xvafund/firm_xvafund.py}{26}{34}{Ajustements pour coût et bénéfice de financement, calculés à partir des profils d'exposition actualisés et de la survie conjointe.}\label{lst:rc:funding-margin-and-capital-adjustments:fva}

\section{Marge initiale et MVA}

Les règles de marge pour les dérivés non compensés (livre~2, chapitre~10) et les chambres de compensation exigent
une marge initiale, versée par les deux parties et ségrégée~: elle ne peut pas être réutilisée, et celui qui la verse
n'en tire qu'un faible rendement tout en la finançant à son propre taux.

\begin{definition}[Ajustement de valeur pour marge]\label{def:rc:funding-margin-and-capital-adjustments:mva}
L'\emph{ajustement de valeur pour marge}\index{ajustement de valeur pour marge} (margin valuation adjustment, MVA) est la
valeur du coût de financement de la marge initiale qu'une position exige sur sa durée de vie~:
$\mathrm{MVA} = \int_0^T(s_f-r_{\mathrm{IM}})\,\E[\mathrm{IM}(t)]\,P(0,t)\,dt$.
\end{definition}

\begin{example}[MVA du swap]\label{ex:rc:funding-margin-and-capital-adjustments:mva}
Avec le modèle de sensibilité du livre~2, la marge initiale vaut 2.326 écarts types de la valeur du
swap sur dix jours, pour une volatilité des taux de 7 points de base par jour~: 4.22~millions de dollars aujourd'hui,
qui décroît à mesure que le DV01 du swap s'éteint (voir la \cref{fig:rc:funding-margin-and-capital-adjustments:im}).
Son financement à 80 points de base coûte 175\,217~dollars sur la durée de vie du swap.
\end{example}

\begin{omfigure}
\begin{tikzpicture}
\begin{axis}[width=11.5cm,height=5.2cm,
  xlabel={temps (années)},ylabel={marge initiale (millions de dollars)},
  tick label style={font=\small},label style={font=\small},
  xmin=0,xmax=10,ymin=0,ymax=4.6,grid=major,grid style={gray!25},
  table/col sep=comma]
\addplot[omDef,very thick,const plot] table[x=t,y=im]{figdata/rates-credit-risk/19-funding-margin-and-capital-adjustments/im.csv};
\end{axis}
\end{tikzpicture}

\omcaption{Marge initiale attendue du swap à dix ans selon un modèle fondé sur la sensibilité~: elle
part de 4.22~millions de dollars et décroît avec le DV01 restant du swap. Son coût de financement sur la
durée de vie est la MVA. Données~: le tutoriel du chapitre.}\label{fig:rc:funding-margin-and-capital-adjustments:im}
\end{omfigure}

\section{Capital~: capital de contrepartie et capital de CVA, et KVA}

Une banque détient des fonds propres face au défaut de ses contreparties et face à la volatilité
de son CVA. Les formules réglementaires transforment les \omterm{def:rc:counterparty-exposure:netting}{ensembles de compensation} en expositions, et les expositions en fonds propres.

\begin{definition}[Exposition en cas de défaut, SA-CCR]\label{def:rc:funding-margin-and-capital-adjustments:saccr}
L'\emph{exposition en cas de défaut}\index{exposition en cas de défaut} (EAD) d'un \omterm{def:rc:counterparty-exposure:netting}{ensemble de compensation} est l'exposition
qu'une règle de fonds propres retient au titre du défaut de sa contrepartie. L'\emph{approche standard
pour le risque de crédit de contrepartie}\index{approche standard pour le risque de crédit de contrepartie}
(SA-CCR) pose $\mathrm{EAD} = 1.4\,(\mathrm{RC}+\mathrm{multiplier}\times\mathrm{AddOn})$~: un coût de remplacement
($\max(V-C,0)$, ou $\max(V-C,\mathrm{Th}+\mathrm{MTA}-\mathrm{NICA},0)$ sous marge) majoré des add-ons prudentiels
par classe d'actifs (pour les taux d'intérêt, 0.5~\% du notionnel multiplié par une durée prudentielle et un
facteur de maturité), réduit pour le collatéral excédentaire.
\end{definition}

\omcode{code/firm/xvafund/firm_xvafund.py}{61}{89}{SA-CCR~: l'add-on de taux d'intérêt avec ses tranches de maturité, le multiplicateur et l'exposition en cas de défaut.}\label{lst:rc:funding-margin-and-capital-adjustments:saccr}

\begin{example}[L'EAD du swap]\label{ex:rc:funding-margin-and-capital-adjustments:ead}
Aujourd'hui le swap n'a aucun coût de remplacement~; sa durée prudentielle est de 7.87 ans et son add-on vaut
0.5~\% de 100 millions de dollars multipliés par 7.87, soit 3.93 millions de dollars, de sorte que son EAD est de 5.51 millions de dollars. Sous un
CSA avec une période de marge en risque de dix jours, le facteur de maturité vaut $1.5\sqrt{10/250} = 0.30$ et
le plancher du coût de remplacement, la MTA, vaut 0.5 million de dollars~: EAD de 2.35 millions de dollars. L'implémentation
reproduit l'exemple travaillé du Comité de Bâle (une EAD de 569 milliers sur trois transactions).
\end{example}

\begin{definition}[Exigence de fonds propres au titre du risque de CVA]\label{def:rc:funding-margin-and-capital-adjustments:cvacap}
L'\emph{exigence de fonds propres au titre du risque de CVA}\index{exigence de fonds propres au titre du risque de CVA} est le capital qu'une banque détient face aux pertes
liées aux variations du CVA. Dans l'approche de base réduite, elle vaut $0.65\sqrt{(\rho\sum_c\mathrm{SCVA}_c)^2+(1-\rho^2)\sum_c\mathrm{SCVA}_c^2}$
avec $\rho = 50~\%$ et, par contrepartie, $\mathrm{SCVA} = \mathrm{RW}\cdot M\cdot\mathrm{EAD}\cdot\mathrm{DF}/1.4$~: un poids de risque par secteur et
qualité de crédit (3~\% pour une entreprise industrielle en investment grade), la maturité effective et un
facteur d'actualisation prudentiel.
\end{definition}

\begin{definition}[Ajustement de valeur pour fonds propres, taux de rendement minimal exigé]\label{def:rc:funding-margin-and-capital-adjustments:kva}
L'\emph{ajustement de valeur pour fonds propres}\index{ajustement de valeur pour fonds propres} (KVA) est la valeur
du rendement exigé par les actionnaires sur le capital qu'une transaction consomme sur sa durée de vie,
$\mathrm{KVA} = h\int_0^T\E[K(t)]P(0,t)Q_C(t)Q_B(t)\,dt$, où $h$ est le \emph{taux de rendement
minimal exigé}\index{taux de rendement minimal exigé} (hurdle rate)~: le rendement du capital, au-delà du taux sans risque, que la banque
exige de ses activités.
\end{definition}

\begin{example}[KVA du swap]\label{ex:rc:funding-margin-and-capital-adjustments:kva}
Avec un poids de risque illustratif de 100~\% et un ratio de fonds propres de 8~\%, le capital de contrepartie aujourd'hui est
440\,686~USD~; le capital de CVA, à un poids de risque de 3~\% et une maturité effective de 5.5 ans (le
temps moyen des paiements du swap), soit 368\,944~USD.
Projeté sur la durée de vie du swap, le coût de remplacement étant fixé à l'\omterm{def:rc:counterparty-exposure:profiles}{exposition attendue}, le
capital atteint un maximum de 1.18 million de dollars après 1.25 an
(la \cref{fig:rc:funding-margin-and-capital-adjustments:capital}). À un \omterm{def:rc:funding-margin-and-capital-adjustments:kva}{taux de rendement minimal exigé} de 10~\%, le KVA vaut
558\,718~USD, l'ajustement le plus important de tous.
\end{example}

\begin{omfigure}
\begin{tikzpicture}
\begin{axis}[width=11.5cm,height=5.2cm,
  xlabel={temps (années)},ylabel={capital (milliers de USD)},
  tick label style={font=\small},label style={font=\small},
  xmin=0,xmax=10,ymin=0,ymax=1700,grid=major,grid style={gray!25},
  legend columns=3,legend style={font=\footnotesize,at={(0.5,-0.25)},anchor=north,draw=none,fill=none,/tikz/every even column/.append style={column sep=8pt}},
  table/col sep=comma]
\addplot[omDef,very thick] table[x=t,y=ccr]{figdata/rates-credit-risk/19-funding-margin-and-capital-adjustments/capital.csv};
\addplot[omMeth,very thick] table[x=t,y=cva]{figdata/rates-credit-risk/19-funding-margin-and-capital-adjustments/capital.csv};
\addplot[omThm,very thick,dashed] table[x=t,y=total]{figdata/rates-credit-risk/19-funding-margin-and-capital-adjustments/capital.csv};
\legend{capital de crédit de contrepartie,capital de CVA,total}
\end{axis}
\end{tikzpicture}

\omcaption{Capital projeté du swap non garanti à dix ans~: capital de crédit de contrepartie
issu du SA-CCR et capital de CVA issu de l'approche de base réduite. Le capital augmente à mesure que le
coût de remplacement attendu s'accumule, puis diminue avec la maturité restante. Données~: le
tutoriel du chapitre.}\label{fig:rc:funding-margin-and-capital-adjustments:capital}
\end{omfigure}

\begin{dated}{2026-09}{Capital de contrepartie et capital de CVA}\label{dat:rc:funding-margin-and-capital-adjustments:capital}
L'approche standard du Comité de Bâle pour le \omterm{def:rc:counterparty-exposure:ccr}{risque de crédit de contrepartie} date de mars
2014~; son cadre révisé pour le CVA, avec l'approche de base et l'approche standard, date de juillet 2020,
dans le cadre des normes de Bâle~III finalisées en décembre 2017. La norme du Cadre de Bâle sur le risque de crédit
entre en vigueur le 1 janvier 2023 et une version suivante à compter du 1 janvier
2028. Les agences américaines ont rendu l'approche standard obligatoire pour les banques en approches avancées à partir du
1 janvier 2022~; dans l'Union européenne, elle s'applique depuis le 28 juin 2021 et le cadre révisé pour le CVA depuis
le 1 janvier 2025. Le chapitre~23 suit les dates du risque de marché.
\end{dated}

\section{Ensemble des ajustements}

\begin{example}[Prix du swap]\label{ex:rc:funding-margin-and-capital-adjustments:stack}
Sans collatéral, les ajustements du swap sont~: CVA 272\,251~USD (sur le premier défaut), DVA
210\,351, FCA 127\,400, FBA 185\,797, aucune MVA, KVA 558\,718~; au total 562\,220~USD, soit 6.8
points de base par an sur le taux fixe du swap. Sous un CSA avec marge initiale bilatérale~: CVA 33\,814,
DVA 17\,745, FCA 3\,102, FBA 3\,367, MVA 175\,217, KVA 156\,573~; au total 347\,595~USD, soit 4.2
points de base (la \cref{fig:rc:funding-margin-and-capital-adjustments:stack}). La garantie remplace les coûts de crédit et de financement par des coûts de marge,
et réduit le capital.
\end{example}

\begin{omfigure}
\begin{tikzpicture}
\begin{axis}[width=11.5cm,height=5.4cm,ybar,bar width=10pt,
  symbolic x coords={CVA,-DVA,FCA,-FBA,MVA,KVA,total},xtick=data,xticklabels={CVA,-DVA,FCA,-FBA,MVA,KVA,Total},
  ylabel={milliers de USD},
  tick label style={font=\small},label style={font=\small},
  ymin=-300,ymax=800,grid=major,grid style={gray!25},
  legend columns=2,legend style={font=\footnotesize,at={(0.5,-0.2)},anchor=north,draw=none,fill=none,/tikz/every even column/.append style={column sep=8pt}},
  table/col sep=comma]
\addplot[fill=omThm!60,draw=omThm] table[x=item,y=uncoll]{figdata/rates-credit-risk/19-funding-margin-and-capital-adjustments/stack.csv};
\addplot[fill=omDef!60,draw=omDef] table[x=item,y=csa]{figdata/rates-credit-risk/19-funding-margin-and-capital-adjustments/stack.csv};
\legend{sans CSA,CSA et marge initiale}
\end{axis}
\end{tikzpicture}

\omcaption{Empilement des ajustements du swap à dix ans, avec les signes tels qu'ils entrent dans le prix
facturé au client (la DVA et la FBA le réduisent). Sans collatéral, le KVA et le CVA dominent~;
avec un CSA et une marge initiale, la MVA remplace l'essentiel des coûts de crédit et de financement. Données~:
le tutoriel du chapitre.}\label{fig:rc:funding-margin-and-capital-adjustments:stack}
\end{omfigure}

Les ajustements se chevauchent et interagissent~: la DVA et la FBA monétisent toutes deux le spread de crédit propre de la banque~;
le capital dépend des couvertures de CVA, qui réduisent l'exigence de fonds propres au titre de la CVA mais nécessitent un financement~; la marge réduit la CVA
et la FVA mais crée une MVA. Un desk les calcule sur des scénarios communs et en affiche la somme,
objet du chapitre~20.

\section{Tutoriel~: la pile d'ajustements}\label{tut:rc:funding-margin-and-capital-adjustments:tutorial}

\begin{tutorial}
\textbf{Objectif.} Calculer la FCA, la FBA, la MVA et la KVA du swap à dix ans, avec et sans CSA, et
la pile complète. \textbf{État final~:} les résultats
des \cref{ex:rc:funding-margin-and-capital-adjustments:fva,ex:rc:funding-margin-and-capital-adjustments:mva,ex:rc:funding-margin-and-capital-adjustments:ead,ex:rc:funding-margin-and-capital-adjustments:kva,ex:rc:funding-margin-and-capital-adjustments:stack}
et les graphiques.
\begin{tutsteps}
\item \textbf{Exposition}~: les valeurs du swap du chapitre~17~; décalage~1 pour le défaut, décalage~0 pour le financement.
\item \textbf{Financement et marge}~: \texttt{fca}, \texttt{fba}, \texttt{mva}.
\item \textbf{Capital}~: \texttt{ir\_addon}, \texttt{sa\_ccr\_ead}, \texttt{scva},
\texttt{ba\_cva\_capital}~; \texttt{capital\_profile()}~; \texttt{kva}.
\item \textbf{Pile}~: \texttt{stack\_table()}~; \texttt{fig\_rc\_xvafund.py} écrit les graphiques.
\end{tutsteps}
\textbf{Pour aller plus loin.} Acheter des couvertures de CVA et recalculer l'exigence de fonds propres au titre de la CVA avec l'approche de base
complète~; relever le \omterm{def:rc:funding-margin-and-capital-adjustments:kva}{taux de rendement minimal exigé} à 12~\% et observer l'évolution du prix.
\end{tutorial}

\section{Construction~: financement, marge et fonds propres}\label{bld:rc:funding-margin-and-capital-adjustments:xvafund}

\begin{build}
\bfield{Rôle} Les ajustements de financement, de marge et de fonds propres de la société, ainsi que l'exposition
réglementaire et l'exigence de fonds propres au titre de la CVA qu'ils utilisent~; les entrées des cotations du chapitre~20.
\bfield{Interface} \texttt{fca}, \texttt{fba}~; \texttt{mva}, \texttt{im\_profile} (sur
\texttt{firm\_ccpbasis} du livre~2)~; \texttt{IrTrade}, \texttt{supervisory\_duration},
\texttt{maturity\_factor}, \texttt{ir\_addon}, \texttt{fx\_addon}, \texttt{multiplier},
\texttt{sa\_ccr\_ead}~; \texttt{scva}, \texttt{ba\_cva\_capital}~; \texttt{kva}.
\bfield{Règles} SA-CCR pour les seuls ensembles de couverture de taux d'intérêt et de change~; approche de base réduite pour la CVA~;
pondérations de risque et \omterm{def:rc:funding-margin-and-capital-adjustments:kva}{taux de rendement minimal exigé} en entrée~; profils sur la grille du chapitre~17.
\bfield{Tests d'acceptation} \texttt{code/firm/xvafund/tests/}~: l'exemple~1 du Comité de Bâle
(EAD 569)~; le facteur de maturité avec marge~; les bornes du multiplicateur~; l'exigence de fonds propres au titre de la CVA d'une contrepartie unique~;
la FCA d'un profil plat.
\bfield{Pour aller plus loin} Les add-ons de crédit, d'actions et de matières premières~; l'approche standard de l'exigence de fonds propres au titre de la CVA
à partir des sensibilités de CVA~; le capital projeté sur chaque trajectoire~; une courbe de financement au lieu
d'un spread plat.
\end{build}

\begin{omsources}
\item C.\ Burgard et M.\ Kjaer, ``Partial differential equation representations of
derivatives with bilateral counterparty risk and funding costs'', \emph{Journal of Credit
Risk} 7(3), 2011, 75--93.
\item A.\ Castagna, ``Yes, FVA is a cost for derivatives desks'', SSRN, 2012 (une réponse à Hull
et White).
\item L.\ Andersen, D.\ Duffie et Y.\ Song, ``Funding value adjustments'', \emph{Journal of
Finance} 74(1), 2019, 145--192.
\item A.\ Green, C.\ Kenyon et C.\ Dennis, ``KVA: capital valuation adjustment'', SSRN, 2014.
\item Comité de Bâle sur le contrôle bancaire, \emph{The standardised approach for measuring
counterparty credit risk exposures}, 2014~; \emph{Targeted revisions to the credit valuation
adjustment risk framework}, 2020.
\end{omsources}

\section{Exercices}

\begin{exercise}[$\star$]\label{exo:rc:funding-margin-and-capital-adjustments:1}
Un \omterm{def:rc:counterparty-exposure:netting}{ensemble de compensation} présente une EE actualisée plate de 5~millions de dollars pendant quatre ans et la banque se finance à
60 points de base. En négligeant la survie, quelle est la FCA~?
\end{exercise}

\begin{exercise}[$\star$]\label{exo:rc:funding-margin-and-capital-adjustments:2}
Calculer l'add-on SA-CCR et l'EAD d'un nouveau swap non margé de 50~millions de dollars à cinq ans.
\end{exercise}

\begin{exercise}[$\star$]\label{exo:rc:funding-margin-and-capital-adjustments:3}
Pourquoi la FVA nette du swap du chapitre ressort-elle comme un bénéfice~?
\end{exercise}

\begin{exercise}[$\star\star$]\label{exo:rc:funding-margin-and-capital-adjustments:4}
Pourquoi la FBA et la DVA se chevauchent-elles, et comment les banques évitent-elles de les compter toutes les deux~?
\end{exercise}

\begin{exercise}[$\star\star$]\label{exo:rc:funding-margin-and-capital-adjustments:5}
Pourquoi la KVA est-elle plus grande que la CVA pour le swap non collatéralisé~?
\end{exercise}

\begin{exercise}[$\star\star$]\label{exo:rc:funding-margin-and-capital-adjustments:6}
Sous le CSA, la MVA est le poste le plus important. Pourquoi la banque ne peut-elle pas l'éviter~?
\end{exercise}

\begin{exercise}[$\star\star\star$]\label{exo:rc:funding-margin-and-capital-adjustments:7}
\emph{Programmation.} Calculer la FCA du swap pour des spreads de financement de 40 et 120 points de base,
avec et sans le CSA.
\end{exercise}

\begin{exercise}[$\star\star\star$]\label{exo:rc:funding-margin-and-capital-adjustments:8}
\emph{Trouver la faille.} «~La FVA est un problème de notre trésorerie, pas du desk de trading~; nous tarifons au
taux overnight et laissons la trésorerie s'inquiéter du financement.~»
\end{exercise}

\section{Problème~: la charge de financement}

\begin{problem}[{Problème du week-end --- ce que le CSA vaut pour la banque}]\label{pb:rc:funding-margin-and-capital-adjustments:1}
La banque du chapitre se finance à 80 points de base au-dessus du taux overnight. Le trésorier du client
demande ce que la banque donnerait, en points de base courants, pour que le client signe un CSA
sur son swap à dix ans.

\medskip\noindent\textbf{Partie I --- Financement.}
\begin{enumerate}
\item Donner la FCA et la FBA sans CSA, ainsi que la FVA nette.
\item Donner la FCA seule sous forme de spread courant.
\item Donner la FCA sous le CSA et la part qui en est supprimée.
\item Pourquoi le besoin de financement dépend-il de la couverture plutôt que de la transaction du client~?
\item Laquelle de la FCA et de la FBA une banque qui ne comptabilise que les coûts de financement reporterait-elle~?
\end{enumerate}

\medskip\noindent\textbf{Partie II --- Crédit et capital.}
\begin{enumerate}[resume]
\item Donner la CVA et la DVA avec et sans le CSA.
\item Donner l'EAD aujourd'hui avec et sans le CSA.
\item Donner la KVA avec et sans le CSA.
\item Pourquoi le CSA réduit-il la KVA moins que la CVA~?
\item Qu'apporte la marge initiale~?
\end{enumerate}

\medskip\noindent\textbf{Partie III --- L'offre.}
\begin{enumerate}[resume]
\item Donner l'ajustement total sans et avec le CSA et la marge initiale.
\item Convertir les deux en points de base annuels.
\item Quelle somme la banque pourrait-elle verser au client à la signature, en points de base annuels~?
\item Que devrait déposer le client, et combien cela lui coûterait-il~?
\item Un CSA sans marge initiale serait-il préférable pour les deux parties~?
\end{enumerate}

\medskip\noindent\textbf{Partie IV --- Jugement.}
\begin{enumerate}[resume]
\item Qui gagne le débat sur la FVA pour le pricing de ce swap~?
\item Pourquoi une banque facturerait-elle encore la KVA si le capital n'est pas un coût monétaire~?
\item Comment la qualité de crédit propre de la banque intervient-elle dans chaque ajustement~?
\item Énoncer le \emph{résultat nommé}~: la FCA du swap non collatéralisé à 80 points de base
et la part que le CSA supprime.
\item En une phrase~: que apporte le collatéral à une banque~?
\end{enumerate}
\end{problem}

\section{Questions d'entretien}

\begin{interviewq}[$\star$ \iqroles{trader, banque}]\label{iq:rc:funding-margin-and-capital-adjustments:1}
Expliquez la FVA à partir de la couverture d'un swap non collatéralisé.
\end{interviewq}

\begin{interviewq}[$\star\star$ \iqroles{chercheur}]\label{iq:rc:funding-margin-and-capital-adjustments:2}
Résumez le débat sur la FVA. Quelle est votre position~?
\end{interviewq}

\begin{interviewq}[$\star\star$ \iqroles{risque, banque}]\label{iq:rc:funding-margin-and-capital-adjustments:3}
Détaillez la méthode SA-CCR pour un seul swap de taux d'intérêt.
\end{interviewq}

\begin{interviewq}[$\star\star$ \iqroles{trader}]\label{iq:rc:funding-margin-and-capital-adjustments:4}
Qu'est-ce que la MVA et comment la réduiriez-vous~?
\end{interviewq}

\begin{interviewq}[$\star\star\star$ \iqroles{chercheur, banque}]\label{iq:rc:funding-margin-and-capital-adjustments:5}
Pourquoi la KVA est-elle controversée, et comment la calculeriez-vous pour une nouvelle transaction~?
\end{interviewq}

\begin{interviewq}[$\star\star\star$ \iqroles{développeur}]\label{iq:rc:funding-margin-and-capital-adjustments:6}
Comment calculeriez-vous l'ensemble des \omterm{def:rc:credit-and-debit-valuation-adjustments:xva}{XVA} d'un large portefeuille de façon cohérente et rapide~?
\end{interviewq}
